\begin{solapartat}
L'equació de un MVHS la podem escriure com (també considerem vàlida si enlloc de la funció cos es fa servir la funció sin:
\begin{align*}
&x(t) = A\,\cos(\omega\,t + \phi) \quad \text{\punts{0,1}} \Rightarrow v(t) = -A\,\omega\,\sin(\omega\,t + \phi) \quad \text{\punts{0,1}} \Rightarrow\\
&a(t) = -A\,\omega^2\cos(\omega\,t + \phi) \quad \text{\punts{0,3}} \Rightarrow a_\text{màxima} = A\,\omega^2 = A\,(2\pi\,\nu)^2 \quad \text{\punts{0,3}}\\
&\Rightarrow A = \frac{a_\text{màxima}}{(2\pi\,\nu)^2} = 4,2\cdot 10^{-3}\ \mathrm{m} \quad \text{\punts{0,2}}
\end{align*}
\end{solapartat}

\begin{solapartat}
La constant de recuperació en un MVHS la podem deduir a partir de:
\begin{align*}
&a(t) = -\omega^2\,x(t) \quad \text{\punts{0,2}} \Rightarrow -k\,x = m\,a = -m\,\omega^2\,x \quad \text{\punts{0,2}}\\
&k = m\,\omega^2 \quad \text{\punts{0,2}} = 85\cdot(2\pi\cdot 6)^2 = 1,2\cdot 10^{5}\ \mathrm{N/m} \quad \text{\punts{0,4}}
\end{align*}
\end{solapartat}
